% rubber: shell_escape

\documentclass[a4paper,11pt]{article}
\usepackage[utf8]{inputenc}
\usepackage{polski}
\usepackage[usenames,dvipsnames]{xcolor} % Musi być przed tikzem
\usepackage{pdfpages}
\usepackage{svg}
\graphicspath{{tikz-riichi/}{./}}
\usepackage{tikz-riichi/mahjong-hand}
\usepackage[margin=1in]{geometry}
\usepackage{enumitem}
\usepackage{wrapfig}
\usepackage[colorlinks]{hyperref}
\setlist{nolistsep}
\setlength\parindent{0pt}
\setlength{\parskip}{0.5\baselineskip}
\renewcommand{\familydefault}{\sfdefault}

% Ten sam kolor rewersu co w obrazkach z oryginalnych zasad.
\definecolor{tilebackcolor}{HTML}{4f4c4d}

\begin{document}
    \includepdf{frontpage/frontpage.pdf}
    % Na potrzeby spisu treści zerujemy parskip (żeby się za bardzo nie rozlazło)
    % i nie kolorujemy linków.
    {
        \setlength{\parskip}{0pt}
        \hypersetup{hidelinks}
        \tableofcontents
    }
    \pagebreak
    \section*{Przedmowa}
\addcontentsline{toc}{section}{\protect\numberline{}Przedmowa}

W ostatnich latach nastąpił gwałtowny i entuzjastyczny globalny rozwój mahjonga.
Coraz więcej graczy z wielu krajów zaczyna doceniać walory tej najbardziej fascynującej gry świata.
W Europie liczba turniejów zauważalnie rośnie, tak samo zresztą, jak liczba klubów mahjongowych.

Dzisiejszy zglobalizowany świat z szybkim Internetem pozwala graczom z różnych krajów, a nawet kontynentów, w prosty jak nigdy sposób wspólnie zasiąść do gry przy wirtualnym stole.
Internetowe gry i dyskusje dotyczą wielu zestawów zasad, ale wariant japoński, Riichi, zawsze cieszył się szczególnym statusem oraz wyjątkowo entuzjastycznymi zwolennikami --- prawdopodobnie zainspirowanymi japońskimi profesjonalnymi graczami, jedynymi mahjongowymi ,,zawodowcami'' na świecie.

Pierwsza edycja reguł Riichi według EMA została wydana podczas przygotowań do Pierwszego Europejskiego Turnieju Riichi w 2008 roku w Hanowerze (Niemcy).
Reguły opierały się na tych, które od lat były niezależnie praktykowane w Holandii i Danii.
W 2012 roku nastąpiła ich rewizja w celu przybliżenia ich do aktualnych reguł japońskiego Riichi, zachowując jednak część różnic.

Podczas przygotowań do Pierwszych Międzynarodowych Mistrzostw Riichi w Paryżu w 2014 roku powstał nowy zestaw zasad, napisany przez Sylvaina Malbeca we współpracy z wieloma osobami z zewnątrz, w tym z Profesjonalną Ligą Mahjonga w Japonii.
Reguły te zostały lekko zmienione w 2015 roku i oczywistym stało się, że w świetle nowopowstałego dzieła potrzebna była rewizja dotychczasowych zasad EMA.

Ta książeczka opisuje standardowe zasady Riichi stosowane przez Europejską Ligę Mahjonga po zrewidowaniu w 2016 roku.

\begin{flushright}
    \emph{
        Tina Christensen\\
        Prezes Europejskiej Ligi Mahjonga\\
        kwiecień 2016
    }
\end{flushright}
    \section*{Podziękowania}
\addcontentsline{toc}{section}{\protect\numberline{}Podziękowania}

Chciałabym podziękować Sjefowi Strikowi, Jenn Barr i Benjaminowi Boasowi za nieocenioną pomoc w zrozumieniu istoty różnic między zestawami reguł, a także za bezcenne porady.

W tej najnowszej wersji dziękuję również Sylvainowi Malbecowi za jego wkład w analizę zasad Riichi stosowanych na świecie oraz skomponowanie reguł Międzynarodowych Mistrzostw Riichi.
Za kulisami Mistrzostw miało miejsce wiele cennych dyskusji, w których brałam udział wraz z Sylvainem Malbeciem, Scottem Millerem, Jenn Barr, Gemmą Collinge oraz Garthem Nelsonem.
Jestem bardzo wdzięczna za porady oraz wgląd w świat profesjonalnych rozgrywek w Japonii uzyskany od ostatnich trzech.
Wreszcie dziękuję Komitetowi Reguł Riichi EMA (Gemma Collinge, Sven-Hendrik Gutsche, Ans Hoogland, Simon Naarman, Krzysztof Sasinowski i Alexander Wankmüller) za owocne dyskusje.
\begin{flushright}
    \emph{
        Tina Christensen\\
        Przewodnicząca Komitetu Reguł Riichi EMA\\
        Europejska Liga Mahjonga\\
        kwiecień 2016
    }
\end{flushright}
    \section*{Wykaz zmian}
\addcontentsline{toc}{section}{\protect\numberline{}Wykaz zmian}

Reguły zostały zaktualizowane w celu usunięcia nieścisłości oraz większej zgodności z aktualnymi regułami japońskiego riichi oraz regułami Międzynarodowych Mistrzostw Riichi.

\textbf{Główne zmiany w stosunku do wersji z 2012 roku:}
\begin{itemize}
    \item Usunięto czerwone piątki;
    \item \texttt{Tanyao} może być otwarte;
    \item \texttt{Renho} jest warte mangana, zamiast yakumana;
    \item \texttt{Daisuushii} jest warte yakumana, zamiast podwójnego yakumana;
    \item Ręka z 13 i więcej \texttt{han} jest warta sanbaimana, zamiast yakumana;
    \item Usunięto wymóg zgłoszenia meldunku w ciągu trzech sekund;
    \item Podmiana kamienia meldunkiem (\texttt{kuikae}) nie jest dozwolona;
    \item Tymczasowy \texttt{furiten} kończy się po doborze kamienia lub meldunku;
    \item Zniesiono wymóg posiadania dwóch \texttt{yaku} przy 5+ honbach;
    \item Usunięto reguły anulujące rozdanie w pierwszej turze (\emph{abortive draws});
    \item Usunięto \texttt{nagashi mangan};
    \item \texttt{Uma} została zmieniona na 15,000/5,000/-5,000/-15,000;
    \item \texttt{Chombo} kosztuje gracza 20,000 i jest naliczane po przyznaniu \texttt{uma};
    \item W niektórych przypadkach zostały złagodzone reguły karania graczy.
\end{itemize}

\textbf{Zmiany w stosunku do reguł Międzynarodowych Mistrzostw Riichi 2015:}
\begin{itemize}
    \item Wygrać może więcej niż jedna osoba;
    \item 4 \texttt{han} 30 \texttt{fu} nie jest zaokrąglane do mangana;
    \item Zmieniają się reguły dotyczące czasu; EMA: \texttt{pon} ma pierwszeństwo przed \texttt{chii};
    \item Zastawy pod \texttt{riichi} na końcu gry wędrują do zwycięzcy.
\end{itemize}

Ten dokument udostępniony jest na licencji Creative Commons Attribution-NonCommercial-ShareAlike 4.0 International.

\begin{flushright}
  \emph{
    Europejskie Stowarzyszenie Mahjonga \\
    kwiecień 2016
  }
\end{flushright}
    \newpage
\section{Kamienie}

Prezentujemy tu 34 podstawowe kamienie do mahjonga.
Pełen zestaw kamieni do mahjonga zawiera cztery identyczne sztuki każdego kamienia.

\subsection{Trzy rodziny}

Są trzy rodziny kamieni.
Każda z nich ma kamienie ponumerowane od jednego do dziewięciu:

\begin{center}
    \begin{tabular}{ll}
        \raisebox{1.2em}{\texttt{manzu} (liczby)}  & \DrawHand{m1~2~3~4~5~6~7~8~9}{3em} \\
        \raisebox{1.2em}{\texttt{pinzu} (monety)}  & \DrawHand{p1~2~3~4~5~6~7~8~9}{3em} \\
        \raisebox{1.2em}{\texttt{souzu} (bambusy)} & \DrawHand{s1~2~3~4~5~6~7~8~9}{3em}
    \end{tabular}
\end{center}

Jedynka bambusów jest zwykle dekorowana ptakiem, którego wzór często jest różny w różnych zestawach kamieni.
Jedynki i dziewiątki zwane są kamieniami skrajnymi (\emph{terminalami}).

\subsection{Honory}

Oprócz kamieni z poszczególnych trzech rodzin, w mahjongu występuje siedem kamieni zwanych \emph{honorami}: cztery wiatry i trzy smoki.
Wiatry są pokazane w kolejności \textbf{wschód--południe--zachód--północ}.
Smoki są pokazane w kolejności: \textbf{biały--zielony--czerwony}.
Wzór białego smoka często jest różny w różnych zestawach; zwykle jest przedstawiony jako pusty kamień lub jako niebieska prostokątna ramka.

\begin{center}
    \begin{tabular}{ll}
        \raisebox{1.2em}{\texttt{kazehai} (wiatry)}  & \DrawHand{z1~2~3~4}{3em} \\
        \raisebox{1.2em}{\texttt{sangenpai} (smoki)} & \DrawHand{z5~6~7}{3em}
    \end{tabular}
\end{center}

\subsection{Dodatkowe kamienie}

Przy czterech sztukach każdego z wyżej wymienionych kamieni, zestaw do mahjonga składa się łącznie ze 136 kamieni.
Czasami zestawy do mahjonga zawierają również inne kamienie: kwiaty, pory roku, lub jokery, które nie są jednak używane w riichi mahjongu.

Japońskie zestawy często zawierają czerwone piątki (\texttt{akapai}).
Czerwone piątki są czasem używane w zastępstwie normalnych piątek w taki sposób, że każda rodzina kamieni zawiera jedną czerwoną piątkę i trzy zwykłe piątki.
Każda czerwona piątka w ręce podnosi jej wartość o 1 \texttt{han}.
Czerwone piątki nie są już używane w zasadach riichi EMA.

\subsection{Dodatkowe wyposażenie}

Zestawy do mahjonga często posiadają znaczniki do wskazywania wiatru stołu i liczmany używane do prowadzenia punktacji.
Te ostatnie są używane również jako liczniki (\texttt{honba}) i zastawy do \texttt{riichi}.
Wartości pokazanych poniżej liczmanów to odpowiednio 100, 1,000, 5,000, i 10,000.
Dodatkowo można użyć liczmanów w innym kolorze do wyznaczania wartości 500.
Gracze rozpoczynają z 30,000 punktów.
Zestaw powinien również zawierać co najmniej dwie kości.

\makeatletter
\definecolor{myred}{HTML}{D4213D}
\newcommand\@blackbigdot[1]{\filldraw[black] #1 circle (0.16);}
\newcommand\@redbigdot[1]{\filldraw[myred] #1 circle (0.16);}
\newcommand\@blacksmalldot[1]{\filldraw[black] #1 circle (0.09);}
\newcommand\@redsmalldot[1]{\filldraw[myred] #1 circle (0.09);}
\newcommand\@drawstick[2]{
    \begin{tikzpicture}[x=#1,y=#1]
        \draw[rounded corners=5,thick] (0,0) rectangle (8,1);
        #2
    \end{tikzpicture}
}
\newcommand\StickA[1]{\@drawstick{#1}{
    \@blacksmalldot{(3.4,0.3)}
    \@blacksmalldot{(3.8,0.3)}
    \@blacksmalldot{(4.2,0.3)}
    \@blacksmalldot{(4.6,0.3)}
    \@blacksmalldot{(3.4,0.7)}
    \@blacksmalldot{(3.8,0.7)}
    \@blacksmalldot{(4.2,0.7)}
    \@blacksmalldot{(4.6,0.7)}
}}
\newcommand\StickB[1]{\@drawstick{#1}{
    \@redbigdot{(4,0.5)}
}}
\newcommand\StickC[1]{\@drawstick{#1}{
    \@redbigdot{(4,0.5)}
    \@redsmalldot{(3.6,0.25)}
    \@redsmalldot{(3.6,0.75)}
    \@redsmalldot{(4.4,0.25)}
    \@redsmalldot{(4.4,0.75)}
}}
\newcommand\StickD[1]{\@drawstick{#1}{
    \@redbigdot{(4,0.5)}
    \@redsmalldot{(3.6,0.25)}
    \@redsmalldot{(3.6,0.75)}
    \@redsmalldot{(4.4,0.25)}
    \@redsmalldot{(4.4,0.75)}
    \@redsmalldot{(3.2,0.5)}
    \@redsmalldot{(4.8,0.5)}
    \@blackbigdot{(2.5,0.5)}
    \@blackbigdot{(5.5,0.5)}
}}
\makeatother


\begin{center}
    \begin{tabular}{ll}
        \raisebox{.7em}{100} & \StickA{2em} \\
        \raisebox{.7em}{1,000} & \StickB{2em} \\
        \raisebox{.7em}{5,000} & \StickC{2em} \\
        \raisebox{.7em}{10,000} & \StickD{2em}
    \end{tabular}
\end{center}
    \newpage
\section{Przygotowanie gry}
\addcontentsline{toc}{section}{\protect\numberline{}Przygotowanie gry}

\subsection{Wiatr gracza}
    \newpage
\section{Przebieg gry}

Celem gracza w pojedynczym rozdaniu (nazywanym \emph{ręką}) jest złożenie
ręki.
Celem gry jest zebranie jak największej łącznej liczby punktów z wygrywających
rąk.
W szczególności nie ma znaczenia, ile rąk złoży gracz --- jedynym kryterium
wygranej jest jego końcowy wynik.

\subsection{Fazy gry}
Tura gracza rozpoczyna się w momencie pozyskania przez niego kamienia i kończy
w momencie odrzucenia kamienia.
Podczas zwyczajnej sekwencji tur każdy z graczy rozgrywa swoją turę dokładnie
raz.
Zwyczajna sekwencja zostaje przerwana, jeśli odrzucony kamień zostaje skradziony
lub też jeżeli gracz w trakcie swojej tury zamelduje \texttt{ankan} (zamknięty
\texttt{kan}). Rozdanie kończy się, gdy któryś z graczy złoży rękę (co daje mu
wygraną), lub też gdy zostanie ogłoszony remis.
Podczas pojedynczej rundy wszyscy gracze kolejno grają jako gracze Wschodu.
Kompletna gra składa się z dwóch rund: rundy wschodu i rundy południa.

\subsection{Ręka}
Kompletna ręka w mahjongu składa się z czterech \emph{grup} oraz pary.
Grupy dzielą się na sekwensy (\texttt{shuntsu}), trójki (\texttt{koutsu})
i czwórki (\texttt{kantsu}).
Dodatkowo, kompletna ręka musi mieć przynajmniej jedno \texttt{yaku} (punktowany
układ).
Gracz będący \texttt{furiten} nie może wygrać na kamieniu odrzuconym przez
innego gracza.

Sekwens to trzy kolejne kamienie liczbowe z tej samej rodziny --- sekwensów nie
można tworzyć ze smoków ani wiatrów.
Kamienie 8-9-1 oraz 9-1-2 z tej samej rodziny nie tworzą sekwensu.

Trójka składa się z trzech identycznych kamieni.

Czwórka składa się z czterech identycznych kamieni.

Para składa się z dwóch identycznych kamieni.

\begin{center}
    \begin{tabular}{lc}
        \raisebox{1.2em}{Sekwens (\texttt{shuntsu})}
            & \DrawHand{m456}{3em} \\
        \raisebox{1.2em}{Trójka (\texttt{koutsu})}
            & \DrawHand{p444}{3em} \\
        \raisebox{1.2em}{Czwórka (\texttt{kantsu})}
            & \DrawHand{z1111}{3em}
    \end{tabular}
\end{center}

W Riichi Mahjongu istnieją dwie specjalne ręce, które nie wymagają posiadania
czterech grup oraz pary: \texttt{chiitoitsu} oraz \texttt{kokushi musou}.

\subsection{Tura gracza}
Gracze rozgrywają kolejno swoje tury.
Zaczyna gracz Wschodu, a po nim kolejno gracze przeciwnie do ruchu wskazówek
zegara.

Gracz rozpoczyna swoją turę dobraniem kamienia z muru.
Jednakże, ponieważ gracz Wschodu rozpoczyna rozdanie z czternastoma kamieniami,
podczas swojej pierwszej tury nie dobiera on kamienia.
Jeśli gracz nie może (lub nie chce) zadeklarować wygranej lub \texttt{kan},
kończy on swoją turę odrzucając jeden ze swoich zasłoniętych kamieni.
Podczas swojej pierwszej tury gracz Wschodu powinien przed odrzutem poczekać, aż
pozostali gracze zobaczą i posortują swoje kamienie.

Gracze powinni zwrócić uwagę na to, by odrzucać kamienie nie zasłaniając ich
ręką.
Odrzuty są umieszczane przed graczem wewnątrz muru w porządku od lewej do
prawej, po sześć kamieni w rzędzie --- tak, aby było klarowne, kto (i w jakiej
kolejności) odrzucił które kamienie.

\subsubsection{Pierwszeństwo i czas przy kradzieży kamieni}
% TODO: Przez dowolnego innego?
% https://discord.com/channels/591384413410164765/656637295491022848/850873563707473920
Ostatno odrzucony kamień może zostać przywłaszczony przez dowolnego gracza do
trójki lub czwórki, dopóki następny gracz nie dobierze kolejnego kamienia z muru.
Na ostatnio odrzuconym kamieniu można zadeklarować wygraną, dopóki następny
gracz nie wykona odrzutu --- wyjątkiem jest sytuacja, w której następny gracz
zadeklaruje wygraną przez dobór (\texttt{tsumo}).

Kradzież do trójki lub czwórki może skutkować utratą kolejki przez graczy,
ponieważ gra zaczyna toczyć się dalej zaczynając od gracza kradnącego, a nie
odrzucającego.
Jeżeli któryś z graczy zamelduje wygraną na odrzuconym kamieniu, wszystkie
równoczesne meldunki kradzieży tego kamienia zostają zignorowane.
Kilku graczy może wygrać na tym samym kamieniu jednocześnie.

Gracz, którego tura ma nastąpić bezpośrednio po odrzucie kamienia, może ukraść
ten kamień do sekwensu.
Jeśli gracz nie zdecyduje się na kradzież, rozpoczyna swoją turę zwyczajnie
dobierając kamień z muru.
\begin{figure}[h]
    \centering
    \includesvg[width=0.7\linewidth]{images/table-midgame.svg}
    \caption{Przykładowy stan stołu w trakcie gry.}
    \label{fig:table-midgame}
\end{figure}

Deklaracja wygranej na odrzuconym kamieniu ma pierwszeństwo przed każdą inną
deklaracją dotyczącą tego kamienia.
Kradzież do trójki lub czwórki ma pierwszeństwo przed kradzieżą do sekwensu.
Gracz, który zadeklarował wygraną, nie może zmienić swojej deklaracji.

Gracze nie mają limitu czasowego na zagrania, ale oczekuje się od nich gry w
rozsądnym tempie.
Gracz dobierający kamień zbyt szybko tak, że inni gracze nie mają czasu wykonać
meldunków, lub wielokrotnie zbytnio przedłużający swoją turę może być decyzją
sędziego ukarany za utrudnianie gry.

Jeśli następny gracz dobierze kamień zbyt szybko uniemożliwiając innemu graczowi
wykonanie meldunku, meldunek uznaje się za ważny, a dobrany kamień umieszcza się
z powrotem w murze.

\subsubsection{Kradzież z podmianą (\texttt{kuikae})}
Kradzież z podmianą (\emph{swap-calling}, \texttt{kuikae}) nie jest dozwolona:
nie jest dozwolona kradzież kamienia i natychmiastowy odrzut tego samego
kamienia.
Ponadto nie jest dozwolona kradzież kamienia do sekwensu i odrzut kamienia z
drugiego końca sekwensu.
\begin{center}
    \begin{tabular}{cccc}
        \raisebox{0.7em}{Nie można ukraść\quad}\DrawHand{p1}{2em}
            & \raisebox{0.7em}{do trójki}
            & \DrawHand{p11*1}{2em}\raisebox{0.7em}{, a następnie odrzucić}
            & \DrawHand{p1}{2em}\raisebox{0.7em}{.} \\
        \raisebox{0.7em}{Nie można ukraść\quad}\DrawHand{p1}{2em}
            & \raisebox{0.7em}{do sekwensu}
            & \DrawHand{p1*23}{2em}\raisebox{0.7em}{, a następnie odrzucić}
            & \DrawHand{p1}{2em}\raisebox{0.7em}{ lub }\DrawHand{p4}{2em}\raisebox{0.7em}{.}
    \end{tabular}
\end{center}

\subsubsection{Otwarty sekwens}
Kamień do sekwensu może być skradziony wyłącznie od gracza po lewej.
Kradzież ostatnio odrzuconego kamienia do sekwensu rozpoczyna się od wyraźnego
powiedzenia na głos ,,chow'' lub ,,chii''.
Następnie gracz odsłania pasujące kamienie ze swojej ręki, po czym odrzuca
kamień ze swojej ręki i zabiera skradzony kamień.
Kolejność dwóch ostatnich czynności nie jest istotna: gracz może najpierw zabrać
odrzucony kamień, a dopiero potem wykonać odrzut.

\subsubsection{Otwarta trójka}
Kradzież ostatnio odrzuconego kamienia do trójki rozpoczyna się od wyraźnego
powiedzenia na głos ,,pung'' lub ,,pon''.
Następnie gracz odsłania pasujące kamienie ze swojej ręki, po czym odrzuca
kamień ze swojej ręki i zabiera skradzony kamień.
Kolejność dwóch ostatnich czynności nie jest istotna: gracz może najpierw zabrać
odrzucony kamień, a dopiero potem wykonać odrzut.

\subsubsection{Otwarta czwórka}
Kradzież ostatnio odrzuconego kamienia do czwórki rozpoczyna się od wyraźnego
powiedzenia na głos ,,kong'' lub ,,kan'' i położenia kamienia awersem do góry
razem z pasującymi trzema kamieniami z ręki gracza.
Po odsłonięciu \texttt{kandora}, gracz dobiera kamień zastępczy z martwego muru
i kontynuuje swoją turę tak, jak gdyby dobrał kamień z muru.

Martwy mur zawsze składa się z 14 kamieni, zatem po każdym zameldowaniu czwórki
ostatni kamień muru staje się częścią martego muru.

\subsubsection{Rozszerzenie otwartej trójki do czwórki}
Otwarta trójka może zostać rozszerzona do czwórki podczas tury gracza po tym,
jak dobrał kamień z muru lub kamień zastępczy z martwego muru (tj. nie w turze,
w której ukradł kamień do trójki lub sekwensu).
Gracz musi powiedzieć na głos ,,kong'' lub ,,kan'' i położyć czwarty kamień
równolegle do obróconego kamienia z rozszerzanej trójki.
Następnie następuje odsłonięcie \texttt{kandora} i dobranie kamienia
zastępczego.
Kamień użyty do rozszerzenia trójki liczy się jako odrzut i można na nim
zadeklarować wygraną.
Martwy mur zawsze składa się z 14 kamieni, zatem po każdym zameldowaniu czwórki
ostatni kamień muru staje się częścią martego muru.

\subsubsection{Zamknięta czwórka}
Zamknięta czwórka może być zadeklarowana podczas tury gracza po tym, jak dobrał
kamień z muru lub kamień zastępczy z martwego muru (tj. nie w turze, w której
ukradł kamień do trójki lub sekwensu).
Gracz musi powiedzieć na głos ,,kong'' lub ,,kan'', odsłonić kamienie składające
się na czwórkę, obrócić dwa środkowe kamienie rewersem do góry, odsłonić
\texttt{kandora} i dobrać kamień zastępczy.
Martwy mur zawsze składa się z 14 kamieni, zatem po każdym zameldowaniu czwórki
ostatni kamień muru staje się częścią martego muru.

Jeśli gracz nie miał otwartych grup, po zameldowaniu zamkniętej czwórki jego
ręka pozostaje zamknięta.

Na kamieniach z zamkniętej czwórki nie można zadeklarować wygranej, z wyjątkiem
wygranej na \texttt{kokushi musou}.

Cztery identyczne kamienie uznaje się za czwórkę wyłącznie wtedy, gdy nastąpiła
deklaracja czwórki.

\begin{figure}[h]
    \centering
    \begin{tikzpicture}
    \tikzset{every node/.style={font=\scriptsize}};
    \node[inner sep=0pt] at (0,0)
        {\includesvg[width=.7\linewidth]{images/dead-wall.svg}};
    \node[anchor=west,align=left] (repl) at (-5, -2.2)
        {Kolejne kamienie\\zastępcze};
    \draw[->,thick] (repl.north) -> (-1,-1);
    \node[align=center] (kandora) at (0, -2.2)
        {Kolejne wskaźniki\\\texttt{kandora}};
    \draw[->,thick] (kandora.north) -> (0.7,-0.9);
    \node[anchor=east,align=right] (wall) at (5, -2.2)
        {Kamienie z muru\\dodawane kolejno\\do martwego muru};
    \draw[->,thick] (wall.north) -> (2.75,-1);
\end{tikzpicture}

    \caption{Szczegółowy opis roli kamieni w martwym murze. Liczba na kamieniu
    odpowiada numerowi czwórki deklarowanej w trakcie rozdania.}
    \label{fig:dead-wall}
\end{figure}

\end{document}
